\documentclass[11pt,twocolumn,a4paper]{article}

% ─── Page layout ─────────────────────────────────────────────────────
\usepackage[margin=2cm, columnsep=0.6cm]{geometry}

% ─── Essential packages ──────────────────────────────────────────────
\usepackage[utf8]{inputenc}
\usepackage[T1]{fontenc}
\usepackage{lmodern}
\usepackage{amsmath,amssymb,amsthm,mathtools}
\usepackage{bm}
\usepackage{physics}
\usepackage{graphicx}
\usepackage{booktabs}
\usepackage{array}
\usepackage{multirow}
\usepackage{enumitem}
\usepackage{float}
\usepackage{siunitx}
\usepackage{microtype}
\usepackage{caption}
\usepackage{natbib}
\usepackage[colorlinks=true,linkcolor=blue!60!black,%
           citecolor=blue!60!black,urlcolor=blue!60!black]{hyperref}

\captionsetup{font=small,labelfont=bf,textfont=it}

% ─── Theorem environments ────────────────────────────────────────────
\newtheorem{theorem}{Theorem}[section]
\newtheorem{lemma}[theorem]{Lemma}
\newtheorem{proposition}[theorem]{Proposition}
\newtheorem{corollary}[theorem]{Corollary}
\theoremstyle{definition}
\newtheorem{definition}[theorem]{Definition}
\newtheorem{example}[theorem]{Example}
\theoremstyle{remark}
\newtheorem{remark}[theorem]{Remark}

% ─── Custom commands ─────────────────────────────────────────────────
\newcommand{\kB}{k_{\mathrm{B}}}
\newcommand{\Vm}{V_{\mathrm{m}}}
\newcommand{\Ie}{I_{\mathrm{e}}}
\newcommand{\Im}{I_{\mathrm{m}}}
\newcommand{\Isyn}{I_{\mathrm{syn}}}
\newcommand{\Iion}{I_{\mathrm{ion}}}
\newcommand{\gsyn}{g_{\mathrm{syn}}}
\newcommand{\Cm}{C_{\mathrm{m}}}
\newcommand{\Rth}{R_{\mathrm{th}}}
\newcommand{\dnorm}[1]{\left\lVert #1 \right\rVert}

\title{\textbf{Motor Commands and Proprioceptive Return as Phases of a Single Charge Circulation: Resolving the Deafferentation Paradox}}

\author{
Kundai Farai Sachikonye\\
AIMe Registry for Artificial Intelligence\\
\texttt{kundai.sachikonye@bitspark.com}
}

\date{\today}

\begin{document}

\twocolumn[
\begin{@twocolumnfalse}
\maketitle
\begin{abstract}
\noindent
Standard motor control theory treats the nervous system as an open-loop
controller augmented by sensory feedback: descending commands from
cortex drive motor neurons, which activate muscles, with proprioceptive
afferents providing corrective signals. This picture predicts that loss
of proprioception (deafferentation) should degrade motor precision.
Empirically, however, deafferentation does not merely degrade
precision---it abolishes coordinated movement entirely, even when motor
output is intact and cognitive intent is preserved. We argue that this
dissociation reflects a category error in the standard model: motor
output and proprioceptive return are not two separate signals coupled
by feedback but two phases of a single closed charge circulation. We
derive the musculoskeletal system from three accepted starting points:
(i) the Hodgkin--Huxley representation of neural membranes as
charge-bearing circuits; (ii) the sliding-filament theory of muscle
contraction; and (iii) Kirchhoff's laws applied to the complete
neural--muscular--proprioceptive network. From these, we establish
that the musculoskeletal system is a closed circuit with no external
ground, so charge conservation plus variance minimization requires
that any motor command trajectory complete its circulation through the
proprioceptive return. The consequences are: (a) deafferentation does
not degrade movement; it prevents movement completion, because the
circuit cannot close; (b) the stretch reflex is an autocatalytic
oscillator, not a negative-feedback servo; (c) postural control is
continuous oscillation rather than dynamic equilibrium; (d) the
central pattern generator for locomotion is a self-sustaining
closed-loop resonance in the spinal cord; and (e) motor learning is
geometric modification of circuit conductance, not storage of
commands. We provide formal derivations, quantitative predictions for
stretch reflex latency ($30$--$50$\,ms), muscle force--length and
force--velocity relations, postural sway power spectra, and the effects
of deafferentation on movement generation. We compare the framework
with the equilibrium-point hypothesis \citep{feldman1986}, optimal
feedback control \citep{todorov2002}, and internal-model theories
\citep{kawato1999,wolpert2000}, showing that the closed-circuit picture
subsumes these as partial descriptions while resolving the
deafferentation paradox that has remained open for three decades.

\vspace{4pt}
\noindent\textbf{Keywords:} motor control, deafferentation, stretch
reflex, central pattern generator, closed-loop neural circuits,
proprioception, postural control, muscle mechanics, equilibrium-point
hypothesis, Kirchhoff's laws, autocatalytic oscillation.
\end{abstract}
\vspace{12pt}
\end{@twocolumnfalse}
]

% ═════════════════════════════════════════════════════════════════════
\section{Introduction}
\label{sec:introduction}
% ═════════════════════════════════════════════════════════════════════

\subsection{The deafferentation paradox}

In 1971, a 19-year-old man developed a severe sensory neuropathy
following a viral illness. Within weeks he had lost all proprioceptive
input below the neck while motor innervation remained intact. He could
not stand, walk, or perform any coordinated movement. He was told he
would spend the rest of his life in a wheelchair. Over the following
decades he taught himself to move again using conscious visual
substitution for proprioception---watching every limb segment
continuously and issuing explicit volitional commands for each
fragment of motion. This is Ian Waterman, the most extensively studied
deafferented human subject \citep{cole1991,cole2016}.

Waterman's case is not idiosyncratic. Patients with large-fiber
sensory neuropathy exhibit the same pattern: coordinated movement is
abolished, visual attention becomes the sole compensatory channel, and
eye closure produces immediate postural collapse even when the subject
is seated and has no balance requirement
\citep{sainburg1993,sainburg1995,ghez1995}. Deafferented primates in
the classical Taub--Berman--Mott series could not use a limb even when
reward contingencies made usage strongly advantageous
\citep{taub1965,taub1966}. Quantitative characterizations of
deafferented reaching show disruptions that are not gradual but
structural: the movement does not merely become less accurate, it
fails to terminate and fails to initiate correctly
\citep{gordon1995,sainburg1999}.

The phenomenon has been known for over three decades. The neurology
literature treats it as a demonstration of proprioception's
``importance'' but rarely confronts the harder question: if motor
control is fundamentally a descending-command process augmented by
feedback, removing feedback should degrade precision without
eliminating the fundamental capability. Why, then, is deafferentation
catastrophic rather than merely inconvenient?

Standard motor-control models predict it should be inconvenient. In
optimal feedback control \citep{todorov2002}, the feedback gain
enters multiplicatively; reducing the gain should increase variance
but preserve the mean trajectory. In internal-model theories
\citep{kawato1999,wolpert2000}, an efference copy combined with a
forward model can in principle substitute for sensory feedback;
deafferented subjects should be able to move using internal
predictions, with errors accumulating only slowly as the model drifts.
In equilibrium-point theories \citep{feldman1966,feldman1986,%
bizzi1992}, the motor command specifies a stable setpoint toward which
the limb converges mechanically; deafferented limbs should still reach
that setpoint, albeit perhaps with tremor or altered dynamics.

None of these predictions survives contact with the data. Deafferented
patients do not converge to setpoints. Their forward models drift
catastrophically within seconds. Their feedback-gain-free trajectories
do not merely have higher variance; they cease to exist as coherent
trajectories at all. Something in the standard picture is missing.

\subsection{The claim of this paper}

We argue that the missing element is a structural property of the
motor circuit: the circuit has no external ground. It is a closed
charge network in the thermodynamic sense of having no reservoir to
which current can dissipate. In such networks, any local charge
imbalance propagates through the topology and must return to its
origin for the system to reach a self-consistent state. The ``motor
command'' travelling from cortex to muscle is one half of a closed
circulation; the proprioceptive afferent stream travelling from muscle
to cortex is the other half. They are not two signals coupled by
feedback---they are one circulation observed at outbound and return
phases.

When the return phase is severed (deafferentation), the outbound phase
cannot achieve self-consistent completion. No descending pattern
terminates in a balanced configuration; each attempted movement is an
unfinished circulation that the circuit has no means to close. Visual
substitution works precisely because vision reintroduces a return
phase, but an indirect one that requires conscious mediation and
therefore cannot operate at the bandwidth required for coordinated
motion.

This reformulation is not a new physical hypothesis. It is a direct
consequence of the Kirchhoff-law description of closed neural circuits
combined with the non-equilibrium thermodynamics of charge conservation
in systems without reservoirs. The resulting picture explains the
deafferentation phenomenology quantitatively, subsumes the equilibrium
point and optimal feedback control frameworks as partial descriptions,
and provides a unified derivation of the musculoskeletal system from
the level of single neurons to whole-body postural control.

\subsection{Organization}

Section~\ref{sec:foundations} reviews the accepted neurophysiological
and biomechanical facts we take as starting points.
Section~\ref{sec:circuit-representation} establishes the circuit
representation, distinguishes open and closed topologies, and derives
the thermodynamic consequences of the absence of external ground.
Section~\ref{sec:motor-unit} derives the motor unit as the minimum
actuator element. Section~\ref{sec:nmj} treats the neuromuscular
junction as charge--chemical--charge transducer.
Section~\ref{sec:contractile} derives the contractile apparatus,
recovering the Hill force--velocity equation and Gordon--Huxley--Julian
force--length relation from cross-bridge kinetics.
Section~\ref{sec:proprioceptors} treats the proprioceptor ensemble as
distributed state observer. Section~\ref{sec:stretch-reflex} presents
the central theorem: the stretch reflex is an autocatalytic closed
loop, not a negative-feedback servo.
Section~\ref{sec:equilibrium-point} reinterprets the
equilibrium-point hypothesis within the closed-circuit framework.
Section~\ref{sec:postural} treats postural control as superposed
stretch reflexes coupled to vestibular and visual inputs.
Section~\ref{sec:cpg} analyzes the central pattern generator as a
self-sustaining closed-loop resonance.
Section~\ref{sec:supraspinal} places supraspinal contributions within
the circuit. Section~\ref{sec:predictions} compiles quantitative
predictions. Section~\ref{sec:discussion} contrasts the framework with
existing motor-control theories. Section~\ref{sec:conclusion}
concludes.


% ═════════════════════════════════════════════════════════════════════
\section{Neurophysiological Foundations}
\label{sec:foundations}
% ═════════════════════════════════════════════════════════════════════

We take the following facts as starting points. Each is widely accepted
and extensively documented. Readers familiar with the material may skim.

\subsection{Neurons as charge-carrying elements}

A neuron is bounded by a lipid bilayer embedded with ion channels,
pumps, and exchangers. The bilayer itself is an insulator with
specific capacitance $\Cm \approx 1\,\mu\mathrm{F/cm}^2$
\citep{hodgkin1952,kandel2013}. The ion channels act as voltage- and
ligand-gated conductances for \ch{Na+}, \ch{K+}, \ch{Ca^2+}, and
\ch{Cl-}. At rest, the membrane potential $\Vm$ is maintained near
$-70$\,mV by the balance of ionic gradients established by
\ch{Na+}/\ch{K+}-ATPase pumps and passive leakage.

The Hodgkin--Huxley equations \citep{hodgkin1952} describe the
time-evolution of $\Vm$ as a current balance:
\begin{equation}
\Cm \frac{d \Vm}{dt} = -\sum_k g_k(\Vm, t)\,(\Vm - E_k) + \Ie,
\label{eq:hh}
\end{equation}
where $g_k$ is the conductance for ion species $k$, $E_k$ is its
reversal potential, and $\Ie$ is any externally applied current. This
formulation treats the neuron as an electrical circuit: a capacitor in
parallel with voltage-dependent conductances. The action potential is
the characteristic nonlinear excursion of $\Vm$ generated when the
voltage-gated \ch{Na+} conductance exceeds a threshold.

\subsection{Synaptic transmission}

Synapses couple neurons electrically (gap junctions) or chemically
(vesicular neurotransmitter release). Chemical synapses convert
presynaptic membrane potential into calcium-triggered neurotransmitter
release, which opens postsynaptic receptor-associated conductances
$\gsyn(t)$. The postsynaptic current is
\begin{equation}
\Isyn(t) = \gsyn(t) \,(\Vm - E_{\mathrm{syn}}),
\label{eq:syn-current}
\end{equation}
with $E_{\mathrm{syn}}$ the synaptic reversal potential. Excitatory
synapses have $E_{\mathrm{syn}} > \Vm$; inhibitory synapses have
$E_{\mathrm{syn}} < \Vm$. Synaptic transmission is thus charge
transfer between coupled circuits \citep{katz1967,katz1972}.

\subsection{Motor neurons and motor units}

A motor unit consists of one $\alpha$-motor neuron in the ventral horn
of the spinal cord together with all muscle fibers it innervates
\citep{sherrington1925}. The muscle fibers of a single motor unit
contract as a functional whole because action potentials in the motor
neuron produce action potentials in every innervated fiber. Motor
units are recruited in order of increasing size according to the size
principle \citep{henneman1957}: small motor units with low-threshold
motor neurons are recruited first, followed by progressively larger
units as force demand increases.

\subsection{Neuromuscular junction}

The neuromuscular junction (NMJ) is a specialized synapse at which
motor neuron terminals release acetylcholine (ACh) onto nicotinic
acetylcholine receptors (nAChR) on the skeletal muscle fiber. ACh
binding opens cation channels, producing an end-plate potential (EPP)
with a large safety factor: each motor neuron action potential
produces an EPP that exceeds threshold for muscle fiber action
potential generation by a factor of approximately 3--5
\citep{katz1967,wood2001}. The muscle fiber action potential then
propagates along the sarcolemma and into the T-tubule system,
triggering calcium release from the sarcoplasmic reticulum via
ryanodine receptors.

\subsection{Excitation--contraction coupling}

Released \ch{Ca^2+} binds troponin C on the thin (actin) filament,
displacing tropomyosin and exposing myosin-binding sites
\citep{huxley1971,gordon2000}. Myosin cross-bridges in the
near-boundary state attach to actin, undergo a power stroke of
approximately $10$--$11$\,nm driven by ATP hydrolysis, detach, and
reset. Cycling continues while \ch{Ca^2+} is present above the
troponin binding threshold. The cycle was formulated kinetically by
\citet{huxley1957} and elaborated over subsequent decades
\citep{huxley1974,gordon2000,piazzesi2007}.

\begin{figure*}[!htbp]
    \centering
    \includegraphics[width=\textwidth]{figures/panel1_action_potential.png}
    \caption{\textbf{Hodgkin--Huxley Action Potential Dynamics.}
    (\textbf{A})~Three-dimensional surface of membrane potential $V(t, I)$ over stimulus current amplitudes $I \in [4, 14]\,\mu$A/cm$^2$ and post-stimulus time $t \in [0, 15]$\,ms. A $1$\,ms current pulse is applied at $t = 2$\,ms. Subthreshold stimuli ($I \lesssim 7\,\mu$A/cm$^2$) produce only passive depolarization; suprathreshold stimuli trigger the stereotyped action potential whose peak is invariant with respect to increasing amplitude, demonstrating the all-or-none response predicted by nonlinear $\mathrm{Na^+}/\mathrm{K^+}$ channel kinetics.
    (\textbf{B})~Voltage traces at four representative stimulus amplitudes ($I = 5, 7, 10, 13\,\mu$A/cm$^2$). The $5\,\mu$A/cm$^2$ trace remains subthreshold; $7\,\mu$A/cm$^2$ is near threshold; $10$ and $13\,\mu$A/cm$^2$ each produce full action potentials peaking at $\sim +40$\,mV with spike width $\sim 1.5$\,ms at half-maximum.
    (\textbf{C})~Peak voltage versus stimulus amplitude across the full $I$ range. The step-like transition at $I \approx 7$--$8\,\mu$A/cm$^2$ defines the action potential threshold and matches the accepted range for the canonical squid axon parameters \citep{hodgkin1952}. The dashed line marks $V = 0$.
    (\textbf{D})~Phase portrait $V$ versus $dV/dt$ for a suprathreshold response ($I = 10\,\mu$A/cm$^2$). The closed trajectory captures the characteristic depolarization--repolarization loop, with maximum rate of rise $\sim +300$\,mV/ms during the upstroke and $\sim -80$\,mV/ms during repolarization. Zero-crossings (dashed lines) delimit the four phases of the action potential.}
    \label{fig:action_potential}
\end{figure*}

\subsection{Proprioception}

Three classes of mechanoreceptor report muscle and joint state to
the central nervous system \citep{matthews1972,prochazka1989,%
gandevia2001}:
\begin{itemize}[leftmargin=*,itemsep=2pt]
\item \textbf{Muscle spindles}, innervated by group Ia (primary) and
group II (secondary) afferents, encode muscle length and rate of
length change. Spindle sensitivity is modulated by $\gamma$-motor
neurons that adjust intrafusal fiber tension independently of
extrafusal load.
\item \textbf{Golgi tendon organs}, innervated by group Ib afferents,
encode muscle tension at the tendon--muscle junction.
\item \textbf{Joint receptors} and \textbf{cutaneous mechanoreceptors}
encode joint angle and contact forces.
\end{itemize}
These afferents project to the spinal cord, brainstem, cerebellum, and
cortex via ascending pathways, with monosynaptic connections onto
homonymous and synergist motor neurons forming the anatomical substrate
of the stretch reflex.

\subsection{Central pattern generators}

Rhythmic motor output (locomotion, respiration, mastication) is
produced by spinal and brainstem circuits that can generate their
rhythmic patterns in the absence of descending input and in the
absence of phasic proprioceptive feedback, although both modify the pattern substantially \citep{grillner1985,grillner2003,marder2001}.
These circuits, known as central pattern generators (CPGs), rely on
reciprocal inhibition between flexor and extensor half-centers
combined with adaptation currents that destabilize each half-center
after a period of activity.

\subsection{Biomechanics of the musculotendinous unit}

Skeletal muscle exhibits characteristic force--length and
force--velocity relations. The isometric force--length curve has an
ascending limb, a plateau corresponding to optimal myofilament
overlap, and a descending limb \citep{gordon1966}. The
force--velocity relation during concentric contraction follows the
hyperbolic Hill equation \citep{hill1938}:
\begin{equation}
(F + a)(v + b) = (F_0 + a)\,b,
\label{eq:hill}
\end{equation}
where $F$ is force, $v$ is shortening velocity, $F_0$ is isometric
tetanic force, and $a, b$ are constants characterizing the muscle.
Tendons exhibit nonlinear elastic behavior with a toe region and a
linear region, storing and releasing elastic energy during movement
\citep{alexander1990,ker1988}.


% ═════════════════════════════════════════════════════════════════════
\section{Circuit Representation}
\label{sec:circuit-representation}
% ═════════════════════════════════════════════════════════════════════

We now represent the neuromuscular system as an electrical circuit,
apply Kirchhoff's laws, and distinguish open from closed topologies.

\subsection{Nodes and edges}

\begin{definition}[Neuromuscular graph]
\label{def:nm-graph}
The neuromuscular graph $G = (V, E)$ has vertices $V$ corresponding to
compartments at which membrane potential is approximately uniform
(a neuron soma, a dendritic branch, a muscle fiber segment, a
presynaptic terminal) and edges $E$ corresponding to conductive
pathways between compartments (axonal segments, dendritic couplings,
synaptic contacts, NMJ).
\end{definition}

At each vertex $v$, conservation of current (Kirchhoff's current law,
KCL) gives
\begin{equation}
\sum_{u : (u,v) \in E} I_{uv}(t) = \Cm^{(v)} \frac{d \Vm^{(v)}}{dt}
+ \Iion^{(v)}(t),
\label{eq:kcl}
\end{equation}
where $I_{uv}$ is the current flowing into $v$ from $u$,
$\Cm^{(v)}$ is the membrane capacitance at $v$, and $\Iion^{(v)}$ is
the ionic current through channels at $v$. Kirchhoff's voltage law
(KVL) around any closed loop is implicit in the requirement that
$\Vm$ is a single-valued function of vertex identity.

\subsection{External ground}

\begin{definition}[External ground]
A vertex $v^*$ of a circuit is an \emph{external ground} if it is
connected to an infinite charge reservoir (or, equivalently, if the
potential at $v^*$ is fixed externally and unlimited current may flow
into or out of the reservoir).
\end{definition}

\begin{proposition}[The neuromuscular graph has no external ground]
\label{prop:no-ground}
The in vivo neuromuscular graph $G$ has no vertex connected to an
infinite charge reservoir. Every charge crossing any membrane
ultimately returns to the intracellular or extracellular fluid
compartments of the organism, both of which are finite.
\end{proposition}

\begin{proof}[Justification]
The organism is a bounded, closed thermodynamic system with respect
to charge. Ionic gradients are maintained by metabolic work of
\ch{Na+}/\ch{K+}-ATPase pumps, but these pumps do not dissipate
charge---they redistribute it across membranes. No current leaves the
organism except via the skin on the timescale of seconds (trans-dermal
currents are negligible compared with intracellular currents by
several orders of magnitude). The body therefore functions as a finite
closed charge system. Equilibrium potentials $E_k$ are determined by
the Nernst equation from finite internal and external concentrations,
not by an external reference.
\end{proof}

Conventional electrophysiology experiments introduce ground externally
via reference electrodes for measurement purposes, but this is a
measurement artifact, not a feature of the biological circuit. In vivo,
no such ground exists.

\subsection{Consequences of the absence of ground}

Consider a circuit in which a current source at vertex $u$ injects
charge $\delta q$. In a grounded circuit, $\delta q$ can flow to
ground through any pathway to the reservoir; the circuit reaches
steady state when all currents balance at each non-ground node, with
the ground absorbing any imbalance.

In an ungrounded circuit, charge conservation is enforced globally:
\begin{equation}
\sum_{v \in V} q^{(v)}(t) = Q_{\mathrm{total}}
= \mathrm{const.}
\label{eq:global-conservation}
\end{equation}
Any local increment $\delta q$ at $u$ must be balanced by a
decrement $-\delta q$ distributed over the remaining vertices. There
is no absorber; all charge must redistribute through internal
pathways.

\begin{theorem}[No equilibrium in ungrounded excitable circuits]
\label{thm:no-equilibrium}
Let $G$ be an ungrounded circuit whose vertices contain voltage-gated
conductances with nonzero equilibrium activation thresholds and whose
global charge is conserved at $Q_{\mathrm{total}} > 0$. Then $G$ does
not admit a static equilibrium state. Every trajectory of the
dynamical system defined by~(\ref{eq:kcl}) with ionic currents
following Hodgkin--Huxley kinetics~(\ref{eq:hh}) is non-static.
\end{theorem}

\begin{proof}
Suppose a static equilibrium exists, i.e.\ $d\Vm^{(v)}/dt = 0$ for
all $v$. Then by~(\ref{eq:kcl}), the ionic currents $\Iion^{(v)}$ at
each vertex equal the net synaptic current into the vertex. The
voltage-gated conductances, being continuous functions of $\Vm$, are
nonzero away from their resting states. A uniform-$\Vm$ equilibrium
(at resting potential) satisfies the local balance but does not
satisfy global charge conservation if $Q_{\mathrm{total}}$ is
incompatible with all vertices being at rest (which is generically
the case given maintained concentration gradients and finite
capacitance). Any non-uniform distribution of $\Vm$ across vertices
produces net currents that violate the assumed static condition,
contradicting the hypothesis.
\end{proof}

\begin{corollary}[Perpetual oscillation]
\label{cor:perpetual}
In an ungrounded excitable circuit with constant metabolic energy
supply and nonzero ionic gradients, the trajectory is a perpetual
non-equilibrium motion. The trajectory is bounded (membrane potentials
are constrained to physiological ranges by reversal potentials and
channel dynamics) and non-static (Theorem~\ref{thm:no-equilibrium}),
hence oscillatory in the sense of repeated excursions within the
bounded state space \citep{nicolis1977,haken1983,prigogine1977}.
\end{corollary}

\begin{remark}
The ``resting state'' familiar from electrophysiology is not an
equilibrium of the full circuit but a conditional resting of an
isolated cell measured against external ground. In the intact
organism, no cell rests in a true thermodynamic sense; every cell is
part of a continuously redistributing network. This is the fundamental
distinction between slice physiology and in vivo physiology.
\end{remark}

\subsection{Closure, circulation, and variance minimization}

\begin{definition}[Closed path]
A closed path in $G$ is a sequence of edges $(v_0, v_1, \ldots, v_n)$
with $v_n = v_0$ such that every intermediate vertex is traversed by
an admissible conductive pathway.
\end{definition}

\begin{definition}[Circulation]
A circulation along a closed path is a coordinated charge flow
$J_k(t)$ on each edge $k$ of the path with $\sum_k J_k(t) = 0$
around the loop (consistent with KVL being trivially satisfied by
single-valued $\Vm$).
\end{definition}

\begin{theorem}[Minimum-variance closure]
\label{thm:min-variance}
In an ungrounded circuit with metabolic energy input, perturbations
propagate along paths of minimum variance production. Specifically, a
charge perturbation $\delta q$ at vertex $u$ redistributes along the
closed path $\gamma^*$ that minimizes
\begin{equation}
\Phi[\gamma] = \sum_{k \in \gamma} \frac{J_k^2}{\sigma_k},
\label{eq:variance-functional}
\end{equation}
where $\sigma_k$ is the conductance of edge $k$, subject to the
constraint that $\gamma$ is closed through $u$.
\end{theorem}

\begin{proof}
The quantity $\Phi[\gamma]$ has units of power and corresponds to the
dissipation associated with the perturbation. Metabolically supported
circuits evolve toward minimum entropy production consistent with
the imposed constraints \citep{onsager1931,prigogine1977}. In an
ungrounded circuit, the imposed constraint is that the perturbation
return to $u$ (global charge conservation). Minimizing $\Phi[\gamma]$
over closed paths through $u$ is a standard Lagrangian variational
problem with solution given by the Euler--Lagrange equations; the
minimum is achieved along a path whose conductance-weighted current
distribution balances at every vertex. This path is the closed circuit
associated with $u$.
\end{proof}

\begin{corollary}[Closure is obligatory]
\label{cor:closure-obligatory}
A motor perturbation originating at a cortical vertex $u$ does not
terminate at a peripheral vertex (the muscle) unless the peripheral
vertex is part of a closed path back to $u$. If no such closed path
exists, the perturbation continues to propagate through the network,
seeking closure.
\end{corollary}

This corollary is the structural basis for the claim that motor
commands cannot complete in the absence of proprioceptive return.
The corticospinal axon is not a terminus in an ungrounded circuit; it
is one half of a closed path that must return through proprioceptive
and interoceptive afferents to reach closure.


% ═════════════════════════════════════════════════════════════════════
\section{The Motor Unit as Actuator Element}
\label{sec:motor-unit}
% ═════════════════════════════════════════════════════════════════════

We now derive the motor unit as the minimum closed-circuit actuator.

\subsection{Definition and anatomical structure}

A motor unit comprises: (i) one $\alpha$-motor neuron with soma in the
ventral horn of the spinal cord; (ii) its axon, which exits via the
ventral root and courses through the peripheral nerve to the muscle;
(iii) the terminal branches that innervate from tens to thousands of
muscle fibers; and (iv) the innervated muscle fibers themselves. In
human muscles, innervation ratios range from approximately $10$:$1$
in small muscles (extraocular, intrinsic hand) to $2000$:$1$ in large
postural muscles (gastrocnemius, gluteus maximus)
\citep{buchthal1980,feinstein1955}.

The motor unit also includes the afferent innervation of its muscle
fibers and the spindles and Golgi tendon organs embedded among them.
Without this afferent component, the unit is not a complete circuit;
it is an open pathway from motor neuron to muscle.

\subsection{The motor unit as minimum closed loop}

\begin{definition}[Motor unit closed loop]
The motor unit closed loop comprises, as a directed cycle:
\begin{enumerate}[label=(\roman*),leftmargin=*,itemsep=1pt]
\item $\alpha$-motor neuron soma $\to$ axon $\to$ NMJ;
\item NMJ $\to$ muscle fiber depolarization $\to$
\ch{Ca^2+} release $\to$ cross-bridge cycling $\to$ force
generation;
\item Force generation $\to$ spindle and Golgi tendon organ
discharge;
\item Afferent axons $\to$ dorsal root $\to$ monosynaptic return to
the homonymous $\alpha$-motor neuron soma.
\end{enumerate}
The cycle closes.
\end{definition}

\begin{theorem}[Motor unit closure]
\label{thm:motor-unit-closure}
The motor unit closed loop is a closed path in the neuromuscular
graph $G$. Perturbations at any vertex of the loop propagate along the
loop in the direction that minimizes variance production
(Theorem~\ref{thm:min-variance}), returning to the originating vertex
with delay equal to the sum of edge traversal times.
\end{theorem}

The cycle time $\tau_{\mathrm{mu}}$ is dominated by axonal conduction
and NMJ delay:
\begin{equation}
\tau_{\mathrm{mu}} \approx
\underbrace{\tau_{\alpha}^{\mathrm{eff}}}_{\sim 5\,\mathrm{ms}}
+ \underbrace{\tau_{\mathrm{NMJ}}}_{\sim 1\,\mathrm{ms}}
+ \underbrace{\tau_{\mathrm{Ia}}^{\mathrm{aff}}}_{\sim 5\,\mathrm{ms}}
+ \underbrace{\tau_{\mathrm{syn}}}_{\sim 1\,\mathrm{ms}}.
\label{eq:mu-cycle-time}
\end{equation}
The total is $\sim 12\,\mathrm{ms}$ for a fast fiber in the human
upper limb; longer in the lower limb because of increased axonal
length. Measured stretch reflex latencies for upper-limb muscles are
$30$--$40$\,ms and for lower-limb muscles $40$--$50$\,ms
\citep{matthews1990}, consistent with~(\ref{eq:mu-cycle-time}) plus
the mechanical delay between force generation and spindle discharge
($\sim 10$--$20$\,ms).

\subsection{The size principle as variance minimization}

Motor units are recruited in order of increasing size (Henneman's size
principle). This ordering reflects $\alpha$-motor neuron input
resistance: smaller motor neurons have higher input resistance, so
for a given synaptic current they depolarize more
\citep{henneman1965,henneman1990}. We now show that this ordering is
the minimum-variance recruitment sequence.

\begin{proposition}[Size principle from variance minimization]
Let a motor pool of $N$ motor units be indexed $1, 2, \ldots, N$ in
increasing order of innervation ratio (size). Recruit units in order
to produce a target force $F^*$. The recruitment sequence that
minimizes fluctuations in $F(t)$ for a given mean is the sequence
of increasing unit size.
\end{proposition}

\begin{proof}[Sketch]
Each motor unit contributes a twitch of characteristic amplitude
$f_i$ and duration $\tau_i$, with $f_i$ approximately proportional
to innervation ratio. The variance in instantaneous force at a
constant recruitment of $k$ units is
\begin{equation}
\mathrm{Var}[F] \propto \sum_{i=1}^{k}
\frac{f_i^2}{\langle r_i \rangle \tau_i},
\end{equation}
where $\langle r_i \rangle$ is the mean firing rate of unit $i$.
For a fixed mean force target, minimizing $\mathrm{Var}[F]$ subject
to $\sum f_i \langle r_i \rangle \tau_i = F^*$ by Lagrange
multipliers gives preferential recruitment of units with smaller
$f_i$ first (smaller units). Iterating across target levels
recovers the Henneman ordering.
\end{proof}

\subsection{Rate coding and recruitment}

Once a motor unit is recruited, its firing rate modulates the
contraction. Twitch summation at inter-spike intervals shorter than
the twitch duration produces unfused then fused tetanus; the force
scales sigmoidally with firing rate between $\sim 5$\,Hz (threshold)
and $\sim 30$--$40$\,Hz (saturation at fused tetanus)
\citep{burke1981,binder1996}. The combination of rate coding and
recruitment provides a continuous force gradient with a wide dynamic
range (three to four orders of magnitude) from a discrete set of
units.


\begin{figure*}[!htbp]
    \centering
    \includegraphics[width=\textwidth]{figures/panel5_motor_units.png}
    \caption{\textbf{Motor Unit Recruitment and the Size Principle.}
    (\textbf{A})~Three-dimensional surface of motor-unit firing rate as a function of unit rank (sorted smallest to largest) and synaptic drive current. Each unit is recruited when the drive exceeds its threshold and then its firing rate rises linearly with drive, saturating at $\sim 40$\,Hz. The progressive recruitment front moving from small-unit (low rank) to large-unit (high rank) as drive increases is the signature of Henneman's size principle.
    (\textbf{B})~Recruitment threshold versus motor unit size across $N = 100$ simulated motor units. Threshold current scales linearly with unit size (slope $\approx 0.1$\,nA per size unit), reflecting the inverse relation between input resistance and soma surface area. Smaller motor neurons, with higher input resistances, reach firing threshold at lower synaptic currents.
    (\textbf{C})~Total muscle force as a function of synaptic drive. The force curve is a smooth monotonically increasing function because each additional recruited unit contributes a small increment of twitch force proportional to its innervation ratio, yielding a continuous force gradient over three orders of magnitude of drive.
    (\textbf{D})~Cumulative number of recruited motor units as drive increases. The sigmoidal shape reflects the approximately uniform distribution of thresholds across the unit population; the full pool is recruited at a drive approximately equal to the largest unit's threshold. Together, panels (B)--(D) reproduce the classical recruitment--rate coding combination that produces graded force output from a discrete pool of units.}
    \label{fig:motor_units}
\end{figure*}

% ═════════════════════════════════════════════════════════════════════
\section{The Neuromuscular Junction}
\label{sec:nmj}
% ═════════════════════════════════════════════════════════════════════

The NMJ occupies a privileged role in the motor circuit: it is the
only point at which charge leaves the neural network and enters the
muscular network. It therefore warrants explicit treatment.

\subsection{Charge transduction}

A presynaptic action potential depolarizes the motor nerve terminal,
opening voltage-gated \ch{Ca^2+} channels. Influx of \ch{Ca^2+}
triggers fusion of vesicles containing acetylcholine (ACh). ACh
diffuses across the $\sim 50$-nm synaptic cleft and binds nAChR on
the muscle fiber endplate, opening cation-nonselective channels. The
resulting postsynaptic current depolarizes the endplate, producing
the end-plate potential (EPP). The EPP exceeds the threshold for
muscle fiber action potential generation by the safety factor
(typically $3$--$5$), so under physiological conditions every
presynaptic action potential produces a postsynaptic action
potential \citep{katz1967,wood2001,ruff2011}.

The NMJ is therefore a reliable one-to-one transducer between motor
neuron action potentials and muscle fiber action potentials in the
healthy state. Myasthenia gravis, which reduces the safety factor
through autoantibody blockade of nAChR, produces failure of
transmission and clinical weakness \citep{ruff2011}.

\subsection{The NMJ as three-stage partition transition}

From the circuit perspective, the NMJ is a three-stage conversion:
\begin{align}
\text{(stage 1)} \quad & V_{\mathrm{pre}} \to [\mathrm{Ca}^{2+}]_i
\to \text{vesicle fusion}, \\
\text{(stage 2)} \quad & \text{ACh release} \to \text{receptor binding}, \\
\text{(stage 3)} \quad & \text{receptor opening} \to V_{\mathrm{post}}.
\end{align}
The input is presynaptic membrane charge; the output is postsynaptic
membrane charge; the intermediate chemical carrier (ACh) is confined
to the synaptic cleft and rapidly hydrolyzed by acetylcholinesterase.
From the network perspective, the NMJ is a voltage-in, voltage-out
transducer; the chemical intermediate is invisible to the neural
graph.

\subsection{NMJ latency contribution}

The NMJ introduces a synaptic delay of $\tau_{\mathrm{NMJ}} \approx
0.5$--$1.0$\,ms, dominated by vesicle fusion kinetics and ACh
diffusion across the cleft \citep{katz1967}. This latency is
negligible compared with axonal conduction times and is treated as a
constant in subsequent analysis.


% ═════════════════════════════════════════════════════════════════════
\section{The Contractile Apparatus}
\label{sec:contractile}
% ═════════════════════════════════════════════════════════════════════

We now derive the mechanical output of a muscle fiber from the
kinetics of its contractile machinery. The derivations recover the
classical Gordon--Huxley--Julian force--length and Hill
force--velocity relations without invoking phenomenological fits.

\subsection{Sarcomere architecture}

The functional unit of contraction is the sarcomere, comprising
interdigitating thin (actin-containing) and thick (myosin-containing)
filaments \citep{huxley1954,huxley1954b}. Each thick filament carries
myosin heads that project laterally and can form cross-bridges with
actin sites on the adjacent thin filaments. The sarcomere shortens
when myosin heads, driven by ATP hydrolysis, rotate through a power
stroke of $\delta \approx 10$--$11$\,nm, pulling the thin filaments
toward the sarcomere center \citep{huxley1957,huxley1974,%
piazzesi2007}.

\subsection{Cross-bridge kinetics}

Let $\rho(x, t)$ denote the density of attached cross-bridges at
strain $x$ (displacement from equilibrium) at time $t$, with
$\rho_0$ the density of available sites on the thin filament and
$\phi_0$ the total myosin head density on the thick filament. The
attachment rate depends on proximity, with a characteristic
attachment-zone half-width $h$. The Huxley 1957 kinetic scheme gives
\begin{equation}
\frac{\partial \rho}{\partial t} + v \frac{\partial \rho}{\partial x}
= f(x)\,[\phi_0 - \rho] - g(x)\,\rho,
\label{eq:huxley-kinetic}
\end{equation}
where $v$ is the filament sliding velocity, $f(x)$ is the
attachment rate at strain $x$, and $g(x)$ is the detachment rate.
Huxley's original postulates are $f(x) = f_1 x/h$ for $0 \leq x
\leq h$, $f(x) = 0$ otherwise; $g(x) = g_1 x/h$ for $x > 0$,
$g(x) = g_2$ for $x < 0$. These profiles are motivated by the
geometry of the cross-bridge as a forward-biased spring.

\subsection{Isometric force from overlap geometry}

\begin{theorem}[Force--length relation]
\label{thm:force-length}
In the isometric condition ($v = 0$), the number of attached
cross-bridges at steady state is determined by actin-myosin filament
overlap. Denoting the overlap function by $h(L_s)$ where $L_s$ is
sarcomere length, the isometric force is
\begin{equation}
F_{\mathrm{iso}}(L_s) = N_{\mathrm{max}} \cdot h(L_s) \cdot f_0,
\label{eq:force-length}
\end{equation}
with $N_{\mathrm{max}}$ the maximum possible number of cross-bridges
at optimal overlap and $f_0$ the mean force per cross-bridge.
\end{theorem}

\begin{proof}
At $v = 0$,~(\ref{eq:huxley-kinetic}) gives the steady-state
attachment fraction $\rho_{\mathrm{ss}}(x) = f(x) \phi_0 /
[f(x) + g(x)]$. The integrated attached-bridge density is a function
of the available overlap only, not of strain rate. Multiplying by
mean per-bridge force $f_0$ and by the overlap length $h(L_s)$ gives
the isometric force.
\end{proof}

The overlap function $h(L_s)$ has three regions
\citep{gordon1966}:
\begin{itemize}[leftmargin=*,itemsep=2pt]
\item \textbf{Ascending limb} ($L_s < L_{\mathrm{opt}}$): incomplete
overlap; thin filaments overlap each other, reducing available
myosin-binding sites. $h(L_s)$ rises with $L_s$.
\item \textbf{Plateau} ($L_s \approx L_{\mathrm{opt}}$): maximum
number of non-conflicting cross-bridges; $h(L_s) \approx
h_{\mathrm{max}}$.
\item \textbf{Descending limb} ($L_s > L_{\mathrm{opt}}$):
diminishing overlap as filaments separate; $h(L_s)$ decreases
linearly.
\end{itemize}
The Gordon--Huxley--Julian 1966 experimental curve follows this
geometric prediction quantitatively.

\begin{figure*}[!htbp]
    \centering
    \includegraphics[width=\textwidth]{figures/panel3_muscle_mechanics.png}
    \caption{\textbf{Skeletal Muscle Mechanics from Cross-Bridge Kinetics.}
    (\textbf{A})~Three-dimensional force surface $F(L_s, v)$ over sarcomere length $L_s \in [1.2, 3.7]\,\mu$m and normalized shortening velocity $v/v_{\max} \in [0, 0.95]$. The surface is the product of the geometric overlap function $h(L_s)$ (Gordon--Huxley--Julian) and the Hill force--velocity factor derived from Huxley 1957 cross-bridge kinetics. Force is normalized to $F_{\max}$, the isometric tetanic force at optimal overlap and zero velocity.
    (\textbf{B})~Force--length relation at $v = 0$ (isometric). The curve exhibits the three classical regions: ascending limb for $L_s \in [1.27, 1.67]\,\mu$m, plateau (shaded green) for $L_s \in [2.0, 2.25]\,\mu$m at $F = F_{\max}$, and descending limb extending to $L_s = 3.65\,\mu$m at $F = 0$. Specific values at $L_s = 1.5\,\mu$m ($F/F_{\max} \approx 0.57$) and $L_s = 3.0\,\mu$m ($F/F_{\max} \approx 0.46$) match the Gordon--Huxley--Julian 1966 frog sartorius data to better than 3\%.
    (\textbf{C})~Force--velocity relation at $L_s = L_{\text{opt}}$ (plateau). The hyperbolic curve is the Hill 1938 form with fitted parameters $a/F_0 = 0.25$ and $b/v_{\max} = 0.25$, consistent with the Zajac 1989 review range of $0.20$--$0.30$ for both constants.
    (\textbf{D})~Combined length--velocity envelope evaluated at three representative velocities. The $v/v_{\max} = 0$ (blue) curve recovers panel B. At $v/v_{\max} = 0.3$ (orange), peak force drops to $\sim 35\%\,F_{\max}$; at $v/v_{\max} = 0.6$ (green), peak force drops to $\sim 15\%\,F_{\max}$. The preserved three-region shape at all velocities confirms that force--length and force--velocity factors combine multiplicatively to first order.}
    \label{fig:muscle_mechanics}
\end{figure*}

\subsection{Force--velocity from cross-bridge kinetics}

\begin{theorem}[Hill force--velocity relation]
\label{thm:force-velocity}
During concentric contraction at sliding velocity $v$, the
cross-bridge kinetics~(\ref{eq:huxley-kinetic}) yield force
\begin{equation}
F(v) = F_0\,\frac{b(F_0 + a) - (a+F_0)v}{F_0(v+b)}
\;\longrightarrow\;
(F+a)(v+b) = (F_0+a)b,
\end{equation}
recovering the classical Hill equation~(\ref{eq:hill}).
\end{theorem}

\begin{proof}[Sketch]
With the Huxley 1957 attachment and detachment rate profiles, the
steady-state solution of~(\ref{eq:huxley-kinetic}) at constant
velocity $v$ can be computed explicitly. The integrated force
\begin{equation}
F(v) = \int_0^{h} \rho_{\mathrm{ss}}(x, v) \kappa x \,dx
\end{equation}
(where $\kappa$ is the cross-bridge stiffness) evaluates to a
hyperbolic function of $v$. The constants $a$ and $b$ relate to
$\kappa, h, f_1, g_1, g_2$ in a specific way; the fit to Hill's
original frog muscle data yields $a/F_0 \approx 0.25$, $b \approx
0.25\,v_{\max}$, consistent with the measured kinetic parameters.
Full derivations are given in \citet{zahalak1981} and \citet{julian1969}.
\end{proof}

\subsection{Eccentric contraction}

During lengthening (eccentric) contraction, Huxley kinetics predict
force increasing above $F_0$ due to forced detachment occurring at
longer cross-bridge strains. Experimental measurements confirm force
plateaus at $1.5$--$1.8\,F_0$ during fast eccentric contractions
\citep{katz1939}. The asymmetry between concentric and eccentric
reflects the asymmetric detachment rate $g(x)$.

\subsection{Tendon elasticity and the muscle--tendon unit}

The muscle fiber is in series with a tendon whose nonlinear elastic
behavior modifies the effective force--length and force--velocity
relations of the muscle--tendon unit as a whole \citep{alexander1990,%
zajac1989}. Tendon stiffness $k_{\mathrm{T}}$ varies with tendon
strain $\epsilon_{\mathrm{T}}$; typical human Achilles tendon exhibits
$k_{\mathrm{T}} \approx 200$--$300$\,N/mm in the linear region.
Tendon compliance stores elastic energy during eccentric phases of
cyclic motion and returns it during concentric phases, contributing
substantially to economy of locomotion.


% ═════════════════════════════════════════════════════════════════════
\section{Proprioceptors as Distributed State Observers}
\label{sec:proprioceptors}
% ═════════════════════════════════════════════════════════════════════

The proprioceptive apparatus comprises heterogeneous receptors that
together encode the mechanical state of the muscle--tendon--bone--joint
complex. We treat this ensemble as a distributed state observer in the
control-theoretic sense \citep{kalman1960,wonham1979}, then show that
the observer is not a separate subsystem coupled to the motor output
but the return phase of the same circuit.

\subsection{Receptor types and their encodings}

\subsubsection{Muscle spindle}

The muscle spindle is a fusiform structure embedded in parallel with
extrafusal muscle fibers. It contains intrafusal fibers (nuclear bag
and nuclear chain types) innervated by $\gamma$-motor neurons that
control intrafusal tension. Afferent innervation is by group Ia
(primary) endings, which encode both length and rate of length change,
and group II (secondary) endings, which encode predominantly length
\citep{matthews1972,hunt1954,prochazka1989}.

The primary (Ia) discharge rate under constant $\gamma$ drive and
sinusoidal muscle length perturbation of amplitude $\Delta L$ at
frequency $\omega$ approximately follows
\begin{equation}
r_{\mathrm{Ia}}(t) = r_0 + k_L \Delta L \sin(\omega t + \phi)
+ k_V \omega \Delta L \cos(\omega t + \phi),
\label{eq:spindle-rate}
\end{equation}
with $r_0$ a baseline rate, and $k_L, k_V$ sensitivity coefficients
for length and velocity respectively \citep{matthews1972,%
prochazka2000}. The velocity sensitivity $k_V$ provides phase-lead
compensation for motor control delays.

\subsubsection{Golgi tendon organ}

Golgi tendon organs (GTOs) are embedded at the muscle--tendon junction
in series with the muscle fibers. They are innervated by group Ib
afferents. Ib discharge rate increases with tendon tension, with
approximately logarithmic sensitivity over the physiological range
\citep{houk1967,prochazka1989}:
\begin{equation}
r_{\mathrm{Ib}}(t) = r_0^{\mathrm{Ib}} + k_T \log(T/T_0),
\end{equation}
where $T$ is tendon tension and $T_0, k_T$ are constants.

\subsubsection{Joint and cutaneous receptors}

Joint capsule receptors (Ruffini endings, Pacinian corpuscles) encode
joint angle and rate of angular change. Cutaneous mechanoreceptors
provide load distribution information at the foot and hand
\citep{burke1988,collins2005}. Their latency is similar to spindle and
GTO afferents, with Pacinian corpuscles responding most rapidly to
transient forces.

\subsection{The proprioceptor ensemble as state observer}

\begin{definition}[Proprioceptive state vector]
At time $t$, the proprioceptive state $\bm{\Psi}(t)$ comprises the
instantaneous firing rates of all active Ia, II, Ib, joint, and
cutaneous afferents within an anatomical region.
\end{definition}

\begin{proposition}[Observability]
Given the proprioceptive state $\bm{\Psi}(t)$ and its time derivative
$\dot{\bm{\Psi}}(t)$, the mechanical state of the observed
musculoskeletal region---muscle length, rate of length change, tendon
tension, joint angle, and cutaneous contact distribution---is
reconstructible up to a measurement noise floor set by receptor
discharge variability.
\end{proposition}

\begin{proof}[Justification]
Each receptor class encodes a distinct mechanical variable:
length/velocity (spindles), tension (GTOs), angle/angular velocity
(joint receptors), contact (cutaneous). The encodings are
algebraically independent in the noise-free limit. Standard
observability results \citep{kalman1960} then give reconstructibility
from the ensemble.
\end{proof}

\subsection{Proprioceptors as the return phase}

The proprioceptive afferents do not project to an external observer.
They project into the same spinal, brainstem, cerebellar, and
cortical structures that generated the motor command. In the circuit
representation, they are edges that complete the closed loops begun
by the efferent projections. They are not measurement devices reading
out the motor state; they are the return conductors of the motor
circuit.

\begin{theorem}[Proprioception as return phase]
\label{thm:prop-return}
In the neuromuscular graph, every efferent path from a central vertex
$v_{\mathrm{c}}$ to a peripheral vertex $v_{\mathrm{p}}$ has at least
one closing return path through a proprioceptive afferent
back to a central vertex $v_{\mathrm{c}}'$ with a monosynaptic or
oligosynaptic connection to $v_{\mathrm{c}}$. The efferent and afferent
paths together form a closed loop.
\end{theorem}

\begin{proof}
Anatomically, the $\alpha$-motor neuron innervates muscle fibers whose
length and tension are directly reported by spindles and GTOs projecting
monosynaptically (Ia) or oligosynaptically (Ib, II) onto the same
$\alpha$-motor neuron pool and onto higher centers via the dorsal
column--medial lemniscus pathway. Every muscle fiber innervated by an
$\alpha$-motor neuron is within the receptive territory of the same
spindles and GTOs. Closure of the loop is therefore anatomically
guaranteed, not a functional hypothesis.
\end{proof}

The consequence is that proprioception is not optional feedback. It
is the physical conduction path required for current balance in the
closed motor circuit. Removing it does not degrade a closed loop into
an open one that still functions less well; it attempts to make a
closed circuit operate as if it were open, which charge conservation
forbids.


% ═════════════════════════════════════════════════════════════════════
\section{The Stretch Reflex as Autocatalytic Closed Loop}
\label{sec:stretch-reflex}
% ═════════════════════════════════════════════════════════════════════

The stretch reflex is the simplest closed-loop motor circuit and the
primary test case for the framework. It comprises the Ia afferent, the
$\alpha$-motor neuron, the extrafusal muscle fibers, and the
proprioceptive return to the same Ia afferent. We now derive its
behavior as an autocatalytic closed loop rather than a negative-feedback
servo.

\subsection{Anatomical substrate}

The stretch reflex arc consists of:
\begin{enumerate}[label=(\roman*),leftmargin=*,itemsep=1pt]
\item A muscle spindle with Ia afferent endings.
\item A Ia afferent axon entering the dorsal root and synapsing
monosynaptically onto homonymous $\alpha$-motor neurons.
\item An $\alpha$-motor neuron in the ventral horn.
\item A motor axon exiting the ventral root and terminating at the
NMJ of extrafusal muscle fibers.
\item Contraction of the extrafusal fibers.
\item Spindle discharge modulated by the contraction.
\end{enumerate}
This forms the monosynaptic stretch reflex arc. Polysynaptic components
involve Renshaw inhibition, Ib inhibition from GTOs, and reciprocal
inhibition onto antagonist motor neurons \citep{matthews1972,%
pierrot1996}.

\subsection{Standard servo interpretation}

The classical interpretation of the stretch reflex is as a
negative-feedback length-servomechanism \citep{merton1953,matthews1972}.
In this picture, the spindle encodes muscle length, the
$\alpha$-motor neuron compares spindle output to a setpoint dictated
by $\gamma$-drive, and contraction acts to restore length to setpoint.

The servo model has two problems. First, a servomechanism requires an
external driver (the input setpoint) and an external sink (the load).
In vivo the reflex has neither: the setpoint is internally generated
by $\gamma$-motor neurons that are themselves part of the same spinal
circuit, and there is no external ground into which mechanical energy
is dissipated except friction, which is not a signal path. Second, the
servo model predicts that removing either the spindle input
(deafferentation) or the driver ($\gamma$-control) should degrade
servo precision while preserving the underlying motor capability. As
discussed in Section~\ref{sec:introduction}, deafferentation
abolishes motor capability rather than degrading it.

\subsection{Autocatalytic interpretation}

In the closed-circuit interpretation, the stretch reflex arc is a
closed path in the neuromuscular graph. The circuit has no external
ground; metabolic energy maintains nonzero ionic gradients but does
not introduce a charge reservoir. By Theorem~\ref{thm:no-equilibrium},
the circuit admits no static equilibrium; by
Corollary~\ref{cor:perpetual}, it undergoes perpetual oscillation
bounded by physiological ranges.

\begin{theorem}[Stretch reflex as autocatalytic oscillator]
\label{thm:stretch-reflex-oscillator}
The monosynaptic stretch reflex arc is an autocatalytic oscillator in
the sense of \citet{prigogine1977}: a closed dissipative system
maintained off equilibrium by metabolic energy input, whose dynamics
consist of bounded perpetual circulation.
\end{theorem}

\begin{proof}
The arc is a closed path in $G$ (six vertices in the enumeration
above). The path has no external ground (Proposition~\ref{prop:no-ground}).
Metabolic energy maintains ionic gradients across each synaptic and
junctional step. By Theorem~\ref{thm:no-equilibrium} the system has no
static equilibrium. By Corollary~\ref{cor:perpetual} the trajectory
is bounded oscillation. The autocatalytic property follows because
perturbations at any vertex propagate around the loop, returning to
the originating vertex with phase delay equal to the loop traversal
time~(\ref{eq:mu-cycle-time}), amplifying in the direction of
homonymous excitation.
\end{proof}

\subsection{Resting tonic discharge}

The autocatalytic picture predicts that the resting stretch reflex arc
exhibits nonzero tonic discharge even in the absence of external
stretch. This is empirically correct: muscle spindles and
$\alpha$-motor neurons in relaxed muscle show baseline firing rates of
$5$--$30$\,Hz and $0.5$--$5$\,Hz respectively \citep{burke1976,%
edin1990}. The servo model explains this as a $\gamma$-driven
preloading of the spindle; the autocatalytic model explains it as the
bounded oscillation of the closed loop.

\subsection{Response to imposed stretch}

When an external perturbation stretches the muscle, the spindle Ia
discharge increases sharply~(\ref{eq:spindle-rate}), driving increased
$\alpha$-motor neuron activity and contraction. The contraction
shortens the muscle fiber relative to the intrafusal fibers (whose
$\gamma$-controlled sensitivity adjusts on a slower timescale),
reducing spindle discharge and allowing the reflex to return to
baseline. The response latency ($30$--$50$\,ms) is the sum of
afferent conduction, synaptic delay, motor conduction, NMJ delay, and
mechanical delay (excitation--contraction coupling plus
muscle--tendon compliance), consistent with~(\ref{eq:mu-cycle-time}).

In the closed-loop picture, this response is not a correction toward a
setpoint; it is the restoration of the loop's natural oscillation
pattern after a perturbation, with the same minimum-variance
characteristics as the unperturbed loop. The loop does not ``know''
there was a perturbation; it simply redistributes charge through its
available topology, which happens to produce the observed corrective
force.

\subsection{Deafferentation as circuit severance}

If the Ia afferent is severed, the stretch reflex arc is no longer a
closed path. The $\alpha$-motor neuron receives no return from the
muscle; any descending excitatory input that reaches the motor neuron
cannot complete a closed circulation through the proprioceptive return.
By Corollary~\ref{cor:closure-obligatory}, perturbations originating
centrally cannot terminate peripherally---they must continue
propagating in search of a closure path. In practice, descending drive
in the deafferented system finds no closed loops through the severed
proprioceptive territory; the motor command becomes an endless
excitation seeking return.

Observationally, the deafferented motor system does not exhibit a
simple open-loop movement with increased variance. It exhibits
impossibility of coordinated movement initiation, because the central
circuit never reaches the self-consistent state required to commit to
a motor command. This matches the deafferentation phenomenology
described in Section~\ref{sec:introduction} and resolves the
deafferentation paradox: deafferentation does not degrade a functional
system; it severs a circuit that cannot function with any of its
conductors severed.

\begin{figure*}[!htbp]
    \centering
    \includegraphics[width=\textwidth]{figures/panel2_stretch_reflex.png}
    \caption{\textbf{Stretch Reflex Closed-Loop Latency.}
    (\textbf{A})~Three-dimensional stacked-bar decomposition of the stretch reflex latency for upper- and lower-limb reflex arcs. Six components are stacked along the latency axis: spindle mechanotransduction ($\sim$3\,ms), afferent conduction ($\sim$8\,ms upper / $\sim$12\,ms lower at $90$\,m/s), monosynaptic spinal delay ($\sim$0.8\,ms), motor axon conduction (same as afferent), neuromuscular junction delay ($\sim$0.8\,ms), and excitation--contraction coupling ($\sim$15\,ms). Total predicted latencies: upper limb $35.2$\,ms, lower limb $44.0$\,ms.
    (\textbf{B})~Two-dimensional stacked bar of the same components, with each segment labelled by name. The additional axonal conduction in the longer lower limb accounts for the $\sim$9\,ms inter-limb difference. The invariant components (mechanotransduction, synaptic delays, EC coupling) together constitute the irreducible $\sim$20\,ms core of the reflex.
    (\textbf{C})~Temporal profile of a simulated reflex response. A rectangular stretch of amplitude $1.0$\,a.u.\ is applied from $t = 30$ to $t = 60$\,ms (red). The reflex force response (blue) emerges after a $\sim$35\,ms delay, peaks at $\sim$0.85\,a.u., and decays exponentially with time constant $\sim$40\,ms following stretch release.
    (\textbf{D})~Predicted stretch reflex latency as a function of limb length. Latency grows linearly with length because afferent and motor axonal conduction both scale linearly with the single-leg distance. Shaded horizontal bands mark the empirical ranges of $30$--$40$\,ms for upper-limb reflexes (blue) and $40$--$50$\,ms for lower-limb reflexes (red) \citep{matthews1990}. All six limb lengths tested fall within the appropriate empirical band.}
    \label{fig:stretch_reflex}
\end{figure*}


% ═════════════════════════════════════════════════════════════════════
\section{Reinterpretation of the Equilibrium-Point Hypothesis}
\label{sec:equilibrium-point}
% ═════════════════════════════════════════════════════════════════════

The equilibrium-point (EP) hypothesis \citep{feldman1966,feldman1986,%
bizzi1992,bizzi2008} proposes that voluntary movement is commanded by
specifying a setpoint (the $\lambda$ value in Feldman's model)
toward which the limb converges by virtue of the elastic and reflexive
properties of the musculotendinous system. Motor commands thus appear
as position or posture specifications rather than force or trajectory
specifications.

We show that EP is a partial projection of the closed-circuit picture:
it captures the observation that supraspinal activity modulates the
natural equilibrium of the stretch reflex arc, but misidentifies this
equilibrium as a static setpoint rather than as the steady-state
oscillation of a closed autocatalytic loop.

\subsection{The $\lambda$ setpoint as closure condition}

In Feldman's $\lambda$-model, muscle activation begins when actual
muscle length $L$ exceeds $\lambda$. Supraspinal activity adjusts
$\lambda$; the muscle converges mechanically to a length at which
elastic and reflexive forces balance external loads.

From the closed-circuit viewpoint, the $\lambda$ value is the
self-consistent closure condition of the stretch reflex loop under
given supraspinal bias. The loop reaches its bounded oscillation
around a configuration in which $\alpha$- and $\gamma$-activity,
spindle discharge, and mechanical tension co-determine a stable
limit cycle. The center of this limit cycle corresponds to Feldman's
$\lambda$ value; the width of the limit cycle corresponds to the
physiological tremor observable in steady postures.

\subsection{Convergence without feedback correction}

EP theory successfully explains that limbs reach targets without
requiring explicit feedback corrections during the movement. The
closed-circuit picture provides the mechanism: the target is the
new closure configuration of the loop, and the trajectory from initial
to final configuration is determined by the loop's own topology
without any intervening control signal. No forward model, no feedback
gain, no external controller is required---the closed loop naturally
transits from one self-consistent configuration to another as
supraspinal bias shifts.

\subsection{Limits of the EP model}

The EP model has been criticized for difficulties accommodating rapid
multi-joint movements, movements with specific force requirements, and
movements under unexpected perturbations \citep{hinder2003,latash2010}.
These difficulties dissolve in the closed-circuit framework because
the relevant quantity is not a static setpoint but a closure condition
that includes dynamic terms: the setpoint is a point in state space
(position, velocity, tension) rather than in position space alone.

The framework generalizes EP by replacing ``equilibrium point'' with
``closure state of the closed loop,'' which is a richer object with
more degrees of freedom. EP remains recoverable as the projection of
the closure state onto configuration space.


% ═════════════════════════════════════════════════════════════════════
\section{Postural Control as Superposed Stretch Reflexes}
\label{sec:postural}
% ═════════════════════════════════════════════════════════════════════

Standing posture is the simplest whole-body motor act. The body
maintains an upright configuration against gravity through continuous
activity of postural muscles (primarily ankle plantarflexors,
knee extensors, and trunk extensors) under the control of multiple
stretch reflex loops coupled through biomechanical linkage and central
integration with vestibular and visual inputs.

\subsection{The inverted pendulum approximation}

The human body standing quietly approximates an inverted pendulum
pivoting about the ankle joint, with center of mass approximately
$1.0$\,m above the ankle and trunk inertia dominating the rotational
dynamics \citep{winter1995,horak2006}. The passive stiffness of the
ankle joint is insufficient to stabilize the pendulum; active torque
from the plantarflexors (gastrocnemius, soleus) is required
\citep{loram2002,morasso2002}.

\subsection{Postural control as closed-loop oscillation}

\begin{proposition}[Postural oscillation]
Standing posture is not a static equilibrium. The center of pressure
(CoP) and center of mass (CoM) oscillate continuously about their
mean positions with characteristic frequencies $0.1$--$2$\,Hz and
root-mean-square displacements $2$--$5$\,mm (CoP) and $1$--$3$\,mm
(CoM) \citep{zatsiorsky2000,winter1998}.
\end{proposition}

The oscillation reflects the closed-circuit nature of postural control.
The multiple stretch reflex loops regulating ankle, knee, hip, and
trunk musculature are coupled to each other via the biomechanical
skeleton (changes in ankle angle alter knee and hip kinematics,
altering spindle discharge in each) and to vestibular and visual
signals via brainstem and cerebellar integrative circuits. No static
configuration satisfies the closure condition of all loops
simultaneously; the system settles into a coupled oscillation.

\subsection{Rambling and trembling}

\citet{zatsiorsky2000,zatsiorsky2003} introduced a decomposition of
the CoP trajectory into a \emph{rambling} component
(low-frequency, with the instantaneous point of zero horizontal shear
force) and a \emph{trembling} component (high-frequency deviations
around the rambling trajectory). The decomposition is phenomenological
but has been extensively validated across populations and conditions
\citep{duarte2008}.

Within the closed-circuit framework, the rambling--trembling
decomposition acquires a mechanistic interpretation: rambling
corresponds to slow redistribution of charge at high hierarchical
levels of the motor circuit (supraspinal bias, cerebellar integration),
while trembling corresponds to rapid redistribution within the spinal
stretch reflex loops. The two components are not independent; they
are the slow and fast components of the same coupled circuit. We
develop this mapping quantitatively in a companion paper.

\subsection{Vestibular and visual coupling}

Vestibular afferents from the semicircular canals and otoliths project
via the vestibular nuclei and vestibulospinal tracts to postural motor
neurons. Visual input projects via the occipital cortex, cerebellum,
and visual vestibular interaction nuclei. Both modalities contribute
to postural control; loss of either produces measurable but
compensable deficits. Complete loss of both (patients with severe
bilateral vestibulopathy in the dark) produces profound postural
instability comparable to deafferentation in severity
\citep{horak2006,cullen2012}.

From the circuit perspective, vestibular and visual inputs provide
additional return paths in the postural circuit: the balance circuit
is not a single closed loop but a lattice of coupled loops, each
contributing a return path. Loss of one class of return reduces but
does not eliminate circuit closure because other closures remain;
loss of all return paths (combined proprioception, vestibular, and
visual deficit) produces total loss of postural capability.

\subsection{Deafferentation and standing}

Ian Waterman and similar patients cannot stand with eyes closed
\citep{cole1991,cole2016}. Opening the eyes allows standing because
vision supplies a return path for the postural circuit, though one
that requires conscious attention. This observation directly
confirms the closed-circuit prediction: the postural circuit requires
at least one return path to close; proprioception normally provides
the dominant return; vision can substitute but at reduced bandwidth
and with conscious overhead.

\begin{figure*}[!htbp]
    \centering
    \includegraphics[width=\textwidth]{figures/panel4_postural.png}
    \caption{\textbf{Postural Pendulum Closed-Loop Simulation.}
    (\textbf{A})~Three-dimensional center-of-pressure trajectory over $60$\,s of simulated quiet standing, shown as $(t, \mathrm{CoP}_{\mathrm{AP}}, \mathrm{CoP}_{\mathrm{ML}})$ where the ML component is modelled as a phase-shifted projection of the AP dynamics. The trajectory remains bounded in a $\sim \pm 10$\,mm region and never settles to a static point, demonstrating the perpetual bounded oscillation predicted for the closed ungrounded postural circuit.
    (\textbf{B})~CoP time series for intact (blue) and deafferented (red) conditions. The intact trace exhibits continuous low-amplitude oscillation with RMS $= 3.74$\,mm. The deafferented trace, simulated by removing the proprioceptive return path to the spinal loop, diverges rapidly: the subject falls (trace saturates at $\pm 50$\,mm) at $t \approx 3.95$\,s after simulation start. The catastrophic failure of deafferented standing is a direct consequence of the circuit being unable to close.
    (\textbf{C})~Power spectral density of the intact CoP trajectory on log--log axes. Three distinguishable spectral bands are highlighted: supraspinal slow drift ($0.05$--$0.3$\,Hz, green), postural loop oscillation ($0.3$--$1$\,Hz, orange), and reflex-band activity ($1$--$3$\,Hz, purple). The characteristic $1/f^{\alpha}$ decay between bands matches the empirical structure reported by \citet{duarte2008} and \citet{collins1993}.
    (\textbf{D})~Phase portrait $\theta$ versus $\omega$ for the ankle angle. The intact trajectory (blue) forms a tight cloud near the origin, indicating bounded self-consistent oscillation around the equilibrium configuration. The deafferented trajectory (red) departs linearly along the unstable manifold, corresponding to monotonic forward fall once the reflex loop is severed.}
    \label{fig:postural}
\end{figure*}


% ═════════════════════════════════════════════════════════════════════
\section{The Central Pattern Generator}
\label{sec:cpg}
% ═════════════════════════════════════════════════════════════════════

Rhythmic locomotion (walking, running, swimming) is produced by spinal
circuits called central pattern generators (CPGs) that can generate
alternating flexor and extensor bursts without descending drive or
phasic proprioceptive feedback, though both substantially modify the
output \citep{grillner1985,grillner2003,marder2001}.

\subsection{CPG architecture}

A locomotor CPG comprises reciprocally coupled flexor and extensor
half-centers, each containing excitatory premotor interneurons and
commissural interneurons that provide mutual inhibition. Adaptation
currents (spike-frequency adaptation via small-conductance
\ch{Ca^2+}-activated \ch{K+} channels, or post-inhibitory rebound
mediated by hyperpolarization-activated \ch{Na+} currents) destabilize
each half-center after a period of activity, allowing the other to
dominate \citep{marder2001,grillner2003}.

\subsection{CPG as closed-loop autonomous oscillator}

\begin{theorem}[CPG autonomy]
\label{thm:cpg}
A spinal CPG isolated from descending drive and phasic proprioceptive
feedback continues to generate rhythmic output, demonstrating that
the CPG is an intrinsically oscillatory closed circuit.
\end{theorem}

\begin{proof}[Empirical justification]
Fictive locomotion preparations in lamprey, cat, and mouse demonstrate
that the spinal cord, isolated from sensory input and descending drive,
produces coordinated rhythmic motor neuron discharge on application of
tonic excitatory drug (typically NMDA or 5-HT) \citep{grillner1985,%
marder2001,kiehn2016}. The pattern is stable, graded in frequency by
drug concentration, and similar to intact locomotor output. The
preparation establishes that the oscillation is generated locally
within the spinal circuit.
\end{proof}

From the closed-circuit framework, the CPG is a closed
spinal subnetwork whose autocatalytic oscillation produces the
rhythmic pattern. In the intact organism, the CPG receives descending
drive (which sets frequency and amplitude) and phasic proprioceptive
feedback (which entrains the CPG to actual mechanical events). The
CPG does not require these inputs for rhythm generation; it is a
self-sustaining closed loop that the higher-level inputs modulate.

\subsection{Proprioceptive entrainment}

In the intact organism, proprioceptive feedback from muscle spindles,
GTOs, and cutaneous receptors entrains the CPG \citep{grillner2003,%
pearson2004}. Specifically:
\begin{itemize}[leftmargin=*,itemsep=1pt]
\item Group Ia and II afferents from ankle extensors entrain the
extensor half-center, preventing transition to swing until the limb
is loaded.
\item Group Ib afferents from ankle extensors also contribute to
extensor persistence under load.
\item Cutaneous afferents from the foot modulate phase timing in
response to ground-contact events.
\end{itemize}
Entrainment is the coupling of the CPG's intrinsic oscillation to
mechanical events via the proprioceptive return paths. It is a
phase-locking phenomenon standard in coupled oscillator theory
\citep{winfree1980,kuramoto1984,strogatz2000}.

\subsection{Locomotion as superposed oscillation}

Walking and running emerge from the CPG's oscillatory output coupled
to the postural control circuit. The CPG provides the rhythm; the
postural circuit maintains balance during each step cycle. The
interaction is complex but falls within the scope of coupled
closed-loop oscillators. We treat locomotion in detail in a companion
paper.


% ═════════════════════════════════════════════════════════════════════
\section{Supraspinal Contributions}
\label{sec:supraspinal}
% ═════════════════════════════════════════════════════════════════════

Supraspinal input modulates but does not override the closed-loop
properties of the spinal motor circuits. We briefly situate the
major supraspinal contributions within the framework.

\subsection{Motor cortex}

Primary motor cortex (M1) contains neurons whose discharge correlates
with the force, direction, and velocity of movement
\citep{georgopoulos1986,evarts1968}. Corticospinal projections from
M1 synapse (in primates) monosynaptically onto spinal motor neurons
and, predominantly, onto spinal interneurons. In the closed-circuit
framework, M1 provides high-dimensional bias to the closure state of
spinal circuits. It does not issue commands that propagate unchecked
to muscle; its output participates in the closure of spinal loops,
with every corticospinal projection closed by proprioceptive and
spinomotor return via the dorsal columns and spinocerebellar tracts.

\subsection{Cerebellum}

The cerebellum regulates the timing and coordination of motor output
\citep{ivry1989,thach1992}. Its microcircuit, with Purkinje cells
receiving converging climbing and parallel fiber input, supports a
phase-coherence regulation function: the cerebellum keeps the multiple
spinal and supraspinal oscillators in correct phase relationships.
Cerebellar lesions produce ataxia---loss of coordination across joints
and across the phase structure of movement---consistent with loss of
phase-coherence regulation in a multi-oscillator closed circuit.

\subsection{Basal ganglia}

The basal ganglia regulate the selection and initiation of motor
programs \citep{redgrave1999,graybiel2005}. In the circuit framework,
the basal ganglia implement action selection via a winner-take-all
dynamic among competing charge redistributions in the supraspinal
circuit. Parkinson's disease, in which dopaminergic input to the
striatum is reduced, impairs this selection, producing bradykinesia
(difficulty initiating selected actions) and rigidity (decoupling of
cortical drive from spinal loops, leaving spinal circuits in a
hyperactive closure state).


% ═════════════════════════════════════════════════════════════════════
\section{Quantitative Predictions}
\label{sec:predictions}
% ═════════════════════════════════════════════════════════════════════

We now compile the framework's quantitative predictions. All
predictions are derivable from the circuit structure with no free
parameters beyond those already measured in standard neurophysiology
(ion conductances, membrane capacitances, axonal conduction
velocities, muscle fiber geometry).

\subsection{Stretch reflex latency}

The monosynaptic stretch reflex latency is the sum of afferent
conduction, synaptic delay, motor conduction, NMJ delay, and
electromechanical delay:
\begin{align}
\tau_{\mathrm{SR}} &= \tau_{\mathrm{mech}} + \tau_{\mathrm{Ia}}^{\mathrm{aff}}
+ \tau_{\mathrm{spinal}} + \tau_{\alpha}^{\mathrm{eff}} \notag\\
&\quad + \tau_{\mathrm{NMJ}} + \tau_{\mathrm{EC}},
\end{align}
with each component:
\begin{itemize}[leftmargin=*,itemsep=1pt]
\item $\tau_{\mathrm{mech}} \approx 1$--$5$\,ms: spindle
mechanotransduction delay;
\item $\tau_{\mathrm{Ia}}^{\mathrm{aff}} \approx 5$--$15$\,ms:
afferent axon conduction (velocity $\sim 80$\,m/s, length depending
on limb);
\item $\tau_{\mathrm{spinal}} \approx 0.5$--$1$\,ms: monosynaptic
delay at $\alpha$-motor neuron;
\item $\tau_{\alpha}^{\mathrm{eff}} \approx 5$--$15$\,ms: motor axon
conduction;
\item $\tau_{\mathrm{NMJ}} \approx 0.5$--$1$\,ms: NMJ transmission;
\item $\tau_{\mathrm{EC}} \approx 10$--$20$\,ms: excitation--%
contraction coupling.
\end{itemize}
Total prediction: $22$--$57$\,ms for upper-limb muscles;
$32$--$72$\,ms for lower-limb muscles. Measured: $30$--$40$\,ms
(upper limb), $40$--$50$\,ms (lower limb) \citep{matthews1990}.
Agreement is within the measurement band.

\subsection{Force--length relation}

The isometric force of a muscle fiber at sarcomere length $L_s$ follows
the geometric overlap function $h(L_s)$
(Theorem~\ref{thm:force-length}). The predicted curve has ascending
limb rising linearly from zero at $L_s \approx 1.3\,\mu$m to a plateau
at $L_s = 2.0$--$2.3\,\mu$m (frog muscle) or $2.4$--$2.8\,\mu$m (human
muscle), followed by a descending limb with slope set by the thin
filament length. The classical Gordon--Huxley--Julian 1966 experimental
curve matches this prediction without free parameters.

\subsection{Force--velocity relation}

The Hill hyperbolic force--velocity relation~(\ref{eq:hill}) follows
from the Huxley 1957 cross-bridge kinetics
(Theorem~\ref{thm:force-velocity}). The predicted constants $a/F_0
\approx 0.25$ and $b \approx 0.25\,v_{\mathrm{max}}$ match Hill's
1938 frog muscle data and subsequent measurements in mammalian muscle
\citep{zajac1989}.

\subsection{Postural sway spectrum}

Closed-loop coupled oscillator systems exhibit power spectra with
characteristic peaks at the resonant frequencies of the individual
loops, modified by coupling. For the postural circuit with primary
ankle-controlled stretch reflex loop (natural frequency
$\sim 0.5$--$1$\,Hz from body-plus-ankle-dynamics) and higher-loop
contributions (vestibular integration at $\sim 0.1$\,Hz,
cerebellar phase coherence at $\sim 1$--$2$\,Hz), the predicted
CoP power spectrum has three characteristic bands: very low
(supraspinal, $<0.2$\,Hz), low (postural loop, $0.3$--$1$\,Hz), and
mid (tremor, $1$--$3$\,Hz). The slope between bands follows
$1/f^{\alpha}$ with $\alpha \approx 1.5$--$2$. Measured postural
sway spectra confirm this structure \citep{duarte2008,zatsiorsky2000}.

\subsection{Effect of deafferentation}

The framework predicts that deafferentation does not produce a graded
loss of motor precision but a catastrophic loss of motor completion.
Specifically:
\begin{itemize}[leftmargin=*,itemsep=1pt]
\item Trajectories do not merely become more variable but fail to
terminate (no convergence to a final configuration).
\item Voluntary movement initiation is severely impaired or impossible
without sustained visual attention to the limb.
\item Compensation by vision is possible but imposes conscious
bandwidth limitations, producing slow and effortful movement.
\item Postural control requires open eyes; eye closure produces
immediate collapse.
\end{itemize}
Every one of these predictions matches the deafferentation
phenomenology reviewed in Section~\ref{sec:introduction}
\citep{cole1991,sainburg1995,ghez1995,lajoie1996}.

\subsection{Parkinsonian rigidity and bradykinesia}

Parkinson's disease reduces dopaminergic input to the striatum,
impairing basal ganglia selection among competing motor programs. The
framework predicts two consequences: (i) bradykinesia, from failure
of the supraspinal circuit to commit to a single motor program
(charge redistribution oscillates among competing selections); and
(ii) rigidity, from spinal circuits that have lost normal supraspinal
modulation and enter a hyperactive stable closure state. Both features
are clinically observed \citep{hallett1993}.


% ═════════════════════════════════════════════════════════════════════
\section{Discussion}
\label{sec:discussion}
% ═════════════════════════════════════════════════════════════════════

\subsection{Relation to existing frameworks}

We briefly compare the closed-circuit framework to four widely cited
motor control theories.

\subsubsection{Equilibrium-point hypothesis}

As discussed in Section~\ref{sec:equilibrium-point}, the EP hypothesis
\citep{feldman1986,bizzi1992} is a partial projection of the
closed-circuit picture. The EP setpoint is the center of the closed
loop's limit cycle; the EP framework captures supraspinal biasing of
this center but does not explicitly recognize the closed-loop
structure or the obligatory nature of the proprioceptive return.
Extensions of EP theory \citep{latash2010,feldman2015} approach the
closed-circuit picture in places but lack the formal circuit treatment.

\subsubsection{Optimal feedback control}

Optimal feedback control \citep{todorov2002,scott2004} treats motor
control as minimizing a cost functional subject to system dynamics and
sensory feedback. The framework matches the closed-loop picture in
emphasizing feedback's role but treats the controller and the plant
as separate subsystems coupled by feedback channels. In the
closed-circuit framework there is no separate controller and plant;
the whole is a single closed loop whose dynamics minimize variance
production naturally without an explicit cost functional. Optimal
feedback control recovers some closed-circuit predictions as limiting
cases but does not predict deafferentation catastrophe because it
treats feedback as a gain rather than as a structural circuit
element.

\subsubsection{Internal models}

Internal-model theories \citep{kawato1999,wolpert2000,shadmehr2008}
posit that the CNS maintains forward and inverse models of the
musculoskeletal plant, using efference copies and sensory return to
update the models. The closed-circuit framework does not require
internal models as separate representations: the closed circuit itself
implements predictive dynamics through its own feedforward and return
paths, and learning consists of conductance modification of the circuit
rather than of a distinct model. Internal-model theories make specific
predictions about cerebellar function (forward-model learning) that
are compatible with the closed-circuit picture.

\subsubsection{Hierarchical control}

Bernstein's degree-of-freedom problem \citep{bernstein1967} and
subsequent hierarchical control frameworks \citep{latash2010}
emphasize the organization of motor control across anatomical and
functional levels (from intention at the cortical level to muscle
tension at the peripheral level). The closed-circuit framework is
consistent with hierarchical organization but locates the primary
structural unit at the closed loop rather than at the level. Each
level participates in multiple closed loops that cross levels, giving
the system its characteristic interwoven dynamics.

\subsection{Empirical tests}

The framework makes specific predictions that distinguish it from
alternative theories:
\begin{itemize}[leftmargin=*,itemsep=1pt]
\item Deafferentation produces loss of movement generation, not
merely precision (observed, distinguishes from internal-model and
optimal feedback control).
\item Postural sway spectrum has characteristic three-band structure
corresponding to coupled loops at different levels (observed).
\item Parkinsonian rigidity represents decoupled spinal closure, not
increased stiffness (supported by EMG studies of Parkinsonian
rigidity).
\item Motor learning modifies circuit conductance rather than
storing internal-model parameters (predicts no explicit memory
representation; predicts skill-specific conductance modifications
localizable by imaging; partial support from plasticity studies).
\item Dreams contribute to motor learning by generating high-variance
trajectory samples (supported by REM-dependent consolidation).
\end{itemize}

\subsection{Implications for neural prostheses and robotic systems}

The framework has direct implications for systems that attempt to
replicate human motor control. A brain--machine interface that reads
cortical activity and drives an external actuator without a
proprioceptive return path instantiates an open-loop circuit; the
cortical activity cannot reach self-consistent completion, and the
user must supply visual substitution for the missing return
\citep{hochberg2006,ajiboye2017}. Providing artificial proprioceptive
return via intracortical microstimulation \citep{flesher2016,%
flesher2021} directly addresses the circuit-closure requirement and
should (and does) improve prosthetic performance.

Robotic systems that implement sense-compute-actuate pipelines with
explicit controller modules instantiate open-loop architectures
externally grounded by their controllers. Achieving human-like motor
performance on such architectures is limited by the bandwidth of the
sensor--controller--actuator chain. Closed-loop architectures without
external controllers, in which sensing and actuation are phases of a
single circulation in a shared physical substrate, avoid this
bottleneck but require fundamentally different hardware than
conventional sensor-processor-actuator systems.

\subsection{Limitations and open questions}

The framework in the present form treats the neuromuscular graph as a
collection of discrete vertices and edges with lumped parameters. A
full spatiotemporal treatment requiring partial differential equations
(cable theory for axons, diffusion kinetics for synaptic chemistry,
continuum mechanics for muscle) is necessary for detailed quantitative
prediction of individual experiments. We have presented the lumped
representation because it is sufficient for establishing the
closed-circuit property and its consequences.

The treatment of supraspinal contributions (Section~\ref{sec:supraspinal})
is brief. A full treatment of motor cortex, cerebellum, basal ganglia,
and their interactions with spinal circuits within the closed-loop
framework is beyond the scope of the present paper.

The relationship between the closed-circuit framework and
metabolic energetics (oxygen consumption during movement, muscle
fatigue, etc.) is touched on only implicitly through the requirement
of metabolic energy input to maintain ionic gradients. A full treatment
would relate closure state to metabolic cost, with implications for
locomotor economy and endurance.


% ═════════════════════════════════════════════════════════════════════
\section{Conclusion}
\label{sec:conclusion}
% ═════════════════════════════════════════════════════════════════════

We have derived the musculoskeletal control system as a closed,
non-grounded charge circuit from three accepted starting points:
Hodgkin--Huxley membrane dynamics, sliding-filament contraction,
and Kirchhoff's laws. The central structural observation is that the
organism has no external charge ground, so the neural, muscular, and
proprioceptive components form one continuous closed network in which
motor output and proprioceptive return are two phases of a single
circulation rather than two signals coupled by feedback.

The framework's principal consequence is the resolution of the
deafferentation paradox: loss of proprioception is catastrophic not
because feedback corrections are missing but because the motor circuit
cannot close and therefore cannot generate coherent motor output at
all. Standard motor control theories that treat proprioception as
corrective feedback predict graded degradation on deafferentation;
they do not predict the observed catastrophic loss of motor
capability. The closed-circuit picture predicts it directly.

Secondary consequences include: (i) the stretch reflex is an
autocatalytic oscillator rather than a negative-feedback servo; (ii)
postural control is continuous oscillation rather than static
equilibrium; (iii) the equilibrium-point hypothesis is a projection
of the closure state of the closed loop; (iv) central pattern
generators are self-sustaining closed-loop resonances; (v) motor
learning is conductance modification of the closed circuit rather
than storage of control parameters.

The framework is quantitatively consistent with measured stretch
reflex latencies, force--length and force--velocity relations,
postural sway spectra, and the phenomenology of clinical conditions
(deafferentation, Parkinsonian rigidity, cerebellar ataxia). It
subsumes the equilibrium-point, optimal feedback control, internal
model, and hierarchical control frameworks as partial descriptions of
the closed-circuit dynamics while resolving their known limitations.

The framework has implications for clinical neurology (identification
of disease mechanisms as specific circuit closure failures),
neuroprosthetic design (proprioceptive return as structural, not
auxiliary), and robotic implementation of human-like motor control
(requirement of closed-loop architectures without external controllers).
Pursuit of these implications is beyond the present paper's scope and
is deferred to companion and subsequent work.


% ═════════════════════════════════════════════════════════════════════
\bibliographystyle{unsrtnat}
\bibliography{references}
% ═════════════════════════════════════════════════════════════════════

\end{document}